% E1 qualitative: residual uncertainty and why the within-surfel term matters
\setlength{\tabcolsep}{1pt}
\begin{tabular}{cccc}
 & \scriptsize $\sqrt{(\Sigma_r)_{uu}}$ & \scriptsize surface part $\sqrt{V_{-j}(u)}$ & \scriptsize tile variance, full / diagonal \\
\rotatebox{90}{\scriptsize\parbox{1.9cm}{\centering weakly supported\\ view 56, 68\% testable}} & \includegraphics[width=0.27\linewidth]{\qfigdir/unc_weak_sd.png} & \includegraphics[width=0.27\linewidth]{\qfigdir/unc_weak_surf.png} & \includegraphics[width=0.27\linewidth]{\qfigdir/unc_weak_ratio.png} \\
\rotatebox{90}{\scriptsize\parbox{1.9cm}{\centering strongly supported\\ view 69, 100\% testable}} & \includegraphics[width=0.27\linewidth]{\qfigdir/unc_strong_sd.png} & \includegraphics[width=0.27\linewidth]{\qfigdir/unc_strong_surf.png} & \includegraphics[width=0.27\linewidth]{\qfigdir/unc_strong_ratio.png} \\
 & \scriptsize 0\,\includegraphics[width=0.11\linewidth,height=0.6em]{\qfigdir/cb_viridis.png}\,7.25 m & \scriptsize 0\,\includegraphics[width=0.11\linewidth,height=0.6em]{\qfigdir/cb_viridis.png}\,7.03 m & \scriptsize 1\,\includegraphics[width=0.11\linewidth,height=0.6em]{\qfigdir/cb_viridis.png}\,9.3 \\
\end{tabular}
